\chapter{Discussion}

\textbf{Status note.} This chapter was originally drafted before the user study (Chapter 6) had
produced any data. Most of its interpretive frame (Sections 7.3–7.4) is kept in that pre-committed
form deliberately. It is the discussion-chapter counterpart of the pre-registered decision table,
fixed before data collection so results cannot be interpreted by a frame built afterwards to flatter
them. The technical benchmark programme of Chapter 5 still has not produced data beyond the offline
detection bake. The user study, however, now has interim results from fourteen Block A pairs and
thirteen Block B pairs (Chapter 6, Section 6.11). Section 7.2 below discusses them against the pre-committed frame that
follows it. This is an interim discussion, not the confirmatory one the frame was built for, and it
is named as such throughout.

\section{What can already be discussed}

\subsection{The design space is itself a result}

The central finding of Part T does not await any benchmark: it is the shape of the design space in
Chapter 3. On a consumer VST headset, diminished reality is not one problem but three, indexed by
compositing access. At Tier 1, the platform grants global colour remapping of the system passthrough
layer (the Styling API's brightness, contrast, saturation, and LUT controls) and nothing spatial: no
per-pixel masks, no world-locked regions, no read access, and, since the deprecation of
surface-projected passthrough at SDK v83, no per-surface control either. The significant fact is
the \textit{absence}: a developer who wants to dim the periphery and spare a window cannot do it at Tier 1
at all. That absence is a platform-policy finding, not an engineering one, and it shapes what DR
research is possible on consumer hardware more than any algorithm does.

The system's response to that absence, the world-locked focus window realised as a shaped \textit{absence}
of overlay, native passthrough showing through it, with all manipulation confined to the periphery,
is an existence proof: windowed subtractive DR is achievable at Tier 2 without OS cooperation, and
hybrid Tier-2/Tier-3 composites (native passthrough inside the window, an effect-processed camera
feed outside it) are achievable with Passthrough Camera Access. The structured prior-art search
reported in Chapter 2, Section 2.5.2 found no publication, product, or open-source project
that composites system passthrough with a live shader-processed camera feed in one view, for any
purpose, while every constituent mechanic has documented precedent. The composite, precisely scoped,
is the novelty claim, and the documented search protocol is what entitles the dissertation to make
it. Should the refresh search before submission surface a concurrent system, the claim degrades
gracefully to independent invention with a characterised design space, which is why the claim is
worded around the space rather than the trick.

\subsection{Engineering lessons that generalise}

Four lessons from the implementation are candidates for reuse beyond this project.

\textbf{Sample by world direction, not screen space, when re-rendering a mono camera feed.}
Screen-space UV mapping of the camera texture gives each eye a different pixel for the same world
point, and the result is double vision. Sampling by world direction through a head-centred pinhole
is eye-independent, so the binocular system fuses the re-rendered image. Any future system that
re-renders a mono VST feed, on this platform or another, will face the same problem and can use
the same solution.

\textbf{Exploit the quality asymmetry: route the task-critical region through the native layer.} The
naive architecture re-renders everything and accepts camera-feed quality everywhere. Inverting it,
manipulating only the periphery and leaving the window as an absence, means the user's
task-relevant vision never pays the Tier-3 quality tax. Chapter 5's T3 benchmark will quantify the
tax. The architectural lesson stands independently: on any platform where an accessible feed is
worse than the native view, put the manipulation where the eyes are not.

\textbf{Decouple effects from head motion.} Head-coupled peripheral restriction increases discomfort
\citep{norouzi2018}. The system therefore fades the effect out during fast head rotation and
restores it on settling, with an immediate restore when the head enters the focus region. This
inverts a common default, effects that follow the head, into one that yields to the head.

\textbf{Consumer-platform DR research inherits platform-policy risk.} The failure museum's least
technical entry may matter most: surface-projected passthrough, a capability this project might have
built on, was deprecated mid-project at SDK v83. Tier boundaries on consumer hardware are set by
vendor policy and move without notice. Research designs that, like this one, keep a working path at
every tier are robust to that movement. Designs premised on a single capability are not.

\subsection{Positioning: software-defined visual noise cancellation}

The literature review's closest conceptual frame for the whole project is visual noise cancellation
(VNC): moderating the visual environment through head-worn hardware the way acoustic noise
cancellation moderates sound. Read through that frame, the present system is a \textit{software-defined}
VNC device: the "cancellation profile" is a shader parameter set, the modes are profiles
(attenuate: Soft/Hard Dark; re-grade: SignPop), and the tier taxonomy describes which profiles the
platform physically permits. The analogy also imports a useful discipline from its acoustic
original: noise-cancelling headphones are evaluated on measured attenuation \textit{and} on what they cost
the wearer in awareness, which is exactly the H1/H2b pairing. The analogy is offered as
positioning, not as a claim of equivalence: acoustic cancellation is subtractive at the signal
level, whereas everything in Chapter 3 exists precisely because consumer VST compositing forbids
signal-level subtraction and forces synthesis by overlay.

\subsection{What could not be discussed before, and what changed}

Before any data existed, no statement about whether the system \textit{works}, in the sense of H1,
H2a, or H2b, was available. The literature gave reason to take both outcomes seriously. Peripheral
distractors reliably impose sustained-attention costs in headsets \citep{ai2025}, and DR-style
attenuation reduced workload in VR simulation \citep{mclaughlin2025}. But area darkening's guidance
effects come mostly from immersive video contexts, desaturation alone has shown weak effects \citep{wang2020}, and awareness costs are documented wherever attenuation succeeds \citep{murph2021}. The design space contained both a working system and a placebo, and only data could say which.
Fourteen Block A pairs and thirteen Block B pairs now exist (Chapter 6, Section 6.11). The next
section discusses what they show, against the interpretive frame that follows it, as an interim reading rather than the
confirmatory one that frame was built for.

\section{Interim results, discussed}

The interim data point toward the working-system end of that range, on both questions the design set
out to ask, with a genuine surprise attached to the second.

\textbf{The workstation benefit (H1) is present and sizeable at this $N$.} There is a 3.34-point
accuracy gain (95\% CI +0.33 to +6.34, $p = .032$, $d_z = 0.641$) on a real-hardware task, with real
physical distractors. It survives both a parametric and an assumption-free test, and ten of fourteen
people improved individually. The honest caveat is the interval's width. A lower bound of +0.33
points is compatible with a benefit barely worth the SESOI the pre-registered design set at 50\% of
the vignette-off error rate. At 60\% power the true effect could plausibly be smaller or larger than
+3.34 once the remaining participants are collected. It is worth noting that adding three
participants moved the estimate by only a third of a point while tightening nothing much, which is
what a real but modest effect looks like as a sample grows.

\textbf{The safety question (H2b) comes back clean.} Pooled hit rate did not drop. It rose by 5.58
points, and the 90\% interval ($-0.98$ to $+12.14$) puts the worst case the data support at about a
one-point loss, far inside the ten-point margin the design set. Section 6.2 committed in advance to
what would refute H2b, a drop greater than ten points alongside a failure to establish equivalence,
and the estimate points the other way, so the refuting observation did not occur. Nothing here
suggests that re-grading salience costs the user their periphery.

\textbf{One methodological lesson comes out of the statistic rather than the result.} Formal
equivalence within $\pm$10 points is not established at this $N$ ($p = .126$), and the reason is
instructive. The two-sided procedure requires ruling out a large \textit{benefit} as well as a large
loss, and the data cannot rule out a benefit above ten points. But a large benefit was never a safety
concern. H2b asks a one-directional question, so the instrument matched to it is a non-inferiority
test against a $-10$-point bound, which is the lower one-sided test and passes at $p = .0006$. The
mismatch is a specification error made when the hypothesis was written, not a property of the system,
and it is worth recording because equivalence testing is increasingly recommended for exactly this
kind of safety claim \citep{lakens2017} and is easy to reach for in its two-sided default form when
the question at hand only has one side. A study writing a safety bound should state the direction it
cares about and test that direction alone.

\textbf{The genuine surprise is where H2a's benefit lives.} The pre-registered hypothesis asked
whether SignPop speeds detection generally. What the interim data show instead is a benefit
concentrated specifically in the 20--30\textdegree{} eccentricity band, the one region where the
colour-pop gate has visible room to act and the target is still close enough to use the result
(+16.5 points, $p = .034$). Every band moves positively, but the other three move by between 1.9 and
5.4 points and none is individually reliable, so one band carries almost all of the signal. This was
not the shape anticipated when the hypothesis was pre-registered. It emerged from breaking the
pooled result down by angle, which makes it a secondary, exploratory finding rather than the
confirmatory test H2a was designed to be.

What makes it more than a lucky slice is that the band was fixed by the design before any data
existed. The soft edge was set at 20\textdegree{} against the useful field of view, the roughly
30\textdegree{} radius within which people extract task-relevant information without moving their
eyes \citep{ball1993}, and set out in Chapter 4, Section 4.2. The benefit therefore appears in precisely the
annulus between where the gate starts having room to act and where the useful field runs out. Inside
20\textdegree{} the gate has nothing to reveal, and beyond 30\textdegree{} it reveals plenty to an eye
that can no longer use it. This is exploratory analysis agreeing with the system's own geometry
rather than fitting noise, which is the strongest kind of exploratory finding available at this
sample size, and it is named as exploratory every time it is used in the rest of this chapter.

\textbf{The subjective and behavioural data partly disagree, and the disagreement is itself
informative.} Perceived focus benefit is strongly positive, 57 of 65 responses. But roughly four in
ten responses on the awareness construct report feeling cut off from the surroundings or noticing
side events late, against a behavioural hit rate that rules out any large peripheral-awareness loss
and, in the 20--30\textdegree{} band, shows an improvement. Two explanations are worth taking seriously rather than
explaining away. First, noticing
that the periphery has been changed is a different experience from missing something in it, and every
DR manipulation studied to date is noticed \citep{sutton2022}. "I can tell my vision has been
altered" and "I failed to detect something" are not the same report, even when only the first is
true. Second, narrowed-field anxiety is documented independently of task performance. It rises in
every restricted-field condition tested in the literature, including conditions where spatial learning
is unaffected \citep{barhorstcates2016}. Reassuring people directly, for instance with a visible
cue when something happens outside the window, would reintroduce exactly the salience the filter
exists to remove. That is the trade-off in its sharpest form, and it is left unresolved here.

\section{The pre-registered interpretation frame}

The decision table (Chapter 6, Section 6.9.3) is reproduced here in condensed form, with the
scientific meaning each row would carry. The dissertation will report the row that occurred once the
final $N = 20$ collection concludes. This section exists so that the meaning of each row is on record
before that data is. The interim data discussed in Section 7.2 sit closest to row 1, but Section 6.9.3
is explicit that this table maps the confirmatory collection, not an interim reading of it.

\textbf{Row 1, H1 supported and H2b passes equivalence}, would validate the mechanism claim on real
hardware. The claim would remain mode-in-context, Hard Dark at a workstation and SignPop on a
monitoring stimulus, and the design cannot support "dark vignettes work" in general. Within that
scope it would be the first performance-based demonstration of DR distraction suppression on consumer
passthrough hardware rather than in VR simulation.

\textbf{Row 2, H1 supported and H2b failing}, is not a failure row. It would reproduce and
\textit{quantify} on real hardware the focus--awareness trade-off documented in simulated DR
\citep{mclaughlin2025,murph2021}, which is arguably the most design-relevant outcome
available, because it converts mode choice from an aesthetic decision into a safety-relevant one:
attenuation modes for closed environments where the periphery genuinely does not matter, salience
modes or nothing where it does.

\textbf{Row 3, H1 refuted}, would be the bounded negative, and it carries a specific scientific
reading. If peripheral visual attenuation at Hard Dark's maximum reduces the error rate by less than
the SESOI, then distractor interference in this paradigm is not carried solely by continuing visual
input. Knowing the tablets are there, and holding an attentional set toward them, may impose costs no
visual overlay can remove. That would be genuinely useful to a field that currently extrapolates from
simulation studies in which removing a distractor is total.

\textbf{Rows 4 and 5} exist to forbid two temptations: retrospectively deciding that a smaller effect
"would also be meaningful", and omitting an abort path from a pre-registered frame.

\section{Block A versus Block B: a built-in test of simulation-based evaluation}

The two confirmatory blocks run the same guidance principle at opposite ends of the feasibility
spectrum: Block A on the real Tier-2 artifact over live passthrough with physical distractors,
Block B on the simulated context, video sphere, scripted probes, that stands in for the
future system. This pairing was designed as a bridge, and it carries an incidental methodological
payload. If the two blocks agree in direction, diminishment helps
focal performance in both, the simulation methodology that the nearest published work relies on
entirely (\citealp{cheng2022}; \citealp{mclaughlin2025}, both VR-simulated) gains a point of external
credibility from this project's real-hardware anchor. If they diverge, the divergence localises
what simulation misses: candidate explanations would include the absence of depth motion in
monoscopic video, the different distractor ontology (baked footage events versus physically present
screens), and hardware-borne effects such as fusion strain that no video condition reproduces.
Either way the comparison is reportable, which is the property that was designed for. The blocks
were not powered as a formal cross-context contrast, and any A-versus-B statement will be framed as
observational.

Operationally, "agreement" will be read at the level of sign and decision-table row, not effect
magnitude: Block A's H1 and Block B's H2a estimate the benefit of diminishment-family manipulations
on focal performance in their respective contexts, and the comparison asks whether both point the
same way. Magnitudes are not comparable across the blocks (different tasks, different dependent
variables, different modes), and no ratio between them will be interpreted.

\section{Transfer: what leaves this dissertation, and under what conditions}

\textbf{To optical see-through (OST) hardware, the dark overlay does not transfer.} The Tier-2 modes
work because a VST display is an opaque screen under full render control. Additive OST optics
cannot darken the world. What transfers is the taxonomy and the question. On OST hardware the
taxonomy collapses toward display-side attenuation, of which segmented dimming hardware (e.g.
Magic Leap 2's segmented dimmer \citep{magic2024}, a dedicated low-resolution panel that locally attenuates light behind masked virtual content) is the nearest existing analogue of a
Tier-2 windowed dim, and toward the additive salience modulation demonstrated by \citet{sutton2022}. Whatever Chapter 6 finds about \textit{what peripheral diminishment does to attention} transfers
as a human-factors result to any hardware that can implement an equivalent stimulus. The
implementation findings transfer only within VST.

\textbf{From the simulation, transfer is conditional, and the conditions are enumerated.} Headset weight
biases only comfort ratings, and conservatively. Attention effects are unaffected in either
direction by lightening hardware. Monoscopic video removes motion-in-depth, so peripheral capture
by looming stimuli is untested. The passenger viewpoint removes task coupling, which is why no
statement in this dissertation concerns driving performance, and none should be quoted as if it
did. Video dynamic range compresses glare, so SignPop's glare-suppression track is exercised only
weakly. Identical stimuli for all participants trade ecological breadth for internal validity, and
the trade is intended. Block A, running on the real system against physical distractors, is the one
component that transfers without these caveats, which is why it carries the primary hypothesis.

\textbf{Oracle conditioning.} Every statement about detector-dependent guidance variants is conditioned
on the oracle assumption, whose measured size is Chapter 5's T4 deliverable. Until the live column
of the oracle-gap table is filled, "the detector sees everything in time" is an assumption with a
plan for its own measurement, not a fact. The study protocol is not insulated from this caveat:
SignPop, Block B's evaluated mode, is itself detector-dependent, on offline, oracle-quality
detection rather than a live pipeline, so RQ2/H2a/H2b inherit the same conditioning as any future
detector-dependent extension, not a lesser version of it.

\section{Limitations, consolidated}

The limitations below are collected from Chapters 3–6 so that they appear once, together, rather
than diluted across the text.

\begin{enumerate}
\item \textbf{Mode $\times$ scenario confound.} Hard Dark is tested only in the workstation context, SignPop only
   in the monitoring context. Every empirical claim is mode-in-context. The design cannot separate
   mode effects from scenario effects, and nothing in the design mitigates it (Chapter 6,
   Section 6.7).
\item \textbf{SignPop's manipulation is a bundle.} A confirmed detection switches on six effects at
   once: colour pop, window carve-out, saturation lift, brightness lift, contrast expansion, and
   darkened surround (Chapter 4, Section 4.5.5). None is isolable in the current design, so
   RQ2/H2a/H2b cannot attribute a measured benefit to any one of them.
\item \textbf{Soft Dark carries no participant data at all.} It was never placed in front of a
   participant (Chapter 6, Section 6.7). Its place in the dissertation rests on design-space
   grounds (Chapter 3, Section 3.5) and on its shader path, never on evidence about how it is
   experienced.
\item \textbf{No eye tracking.} Quest 3 has none. Head pose is a validated proxy at the eccentricities used
   \citep{higgins2022,sitzmann2018}, but covert attention shifts, distraction without
   head movement, are invisible to telemetry. The dependent variables were chosen so that
   \textit{performance outcomes}, not attention allocation, carry the hypotheses. The mechanism remains
   partially inferred.
\item \textbf{Sensitivity floor.} N = 20 gives 80 \% power at dz $\approx$ 0.65 for the classical paired test.
   Trial-level mixed models improve on this, but effects meaningfully smaller than the SESOI are
   detectable only as estimates with wide intervals (decision-table row 4 exists for this reason). At
   the interim sample, achieved power is 60\% for H1 ($N = 14$) and 59\% for the 20--30\textdegree{}
   band ($N = 13$), below the pre-registered target for either. That is why both are reported as interim rather than
   confirmatory findings, regardless of clearing their significance thresholds.
\item \textbf{Single experimenter, single site, single hardware generation.} Scripts and checklists bound
   experimenter variance but cannot eliminate it. Every number is specific to Quest 3 and the
   frozen OS/SDK versions, on a platform whose capabilities have already moved once mid-project.
\item \textbf{Ecological validity boundaries.} The workstation distractors are two tablets with scheduled
   bursts. The monitoring stimulus is a video. The study measures attention mechanisms under
   controlled distraction, deliberately, and its conclusions end where those controls end.
\item \textbf{Residual novelty.} Practice-to-criterion and counterbalancing distribute first-exposure
   effects. They do not abolish the fact that every participant meets peripheral diminishment for
   the first time that hour. Longitudinal adaptation, the deployment-relevant question, is future
   work by construction.
\item \textbf{Passthrough acuity boundary.} The task was designed around the platform's acuity ceiling
   (virtual panel or pilot-verified large print), so no claim extends to fine-detail real-world
   tasks such as reading small text through passthrough, the very context that motivated some of
   the use cases. This is a hardware boundary the dissertation documents rather than solves.
\end{enumerate}

\section{Implications in both futures}

\textbf{If the system works (rows 1–2).} The immediate implication is that subtractive attention
guidance is deliverable \textit{today}, at Tier 2, on unmodified consumer hardware: no OS cooperation, no
eye tracker, no detector. The near-term research programme that opens is concrete: hour-scale field
deployment of the camera-free dark modes (whose endurance T5 tests), the temporal-transition
and user-controlled-release questions already identified as gaps in the literature review (gaps 3
and 4), an eye-tracked replication to convert performance findings into attention-allocation
findings, and an OST translation via segmented dimming hardware. A row-2 outcome adds the design
rule that mode selection is a safety decision, and would motivate a context-detection layer,
which is where the oracle-gap measurement re-enters as an engineering budget.

\textbf{If the system does not work (row 3).} The negative is bounded and useful. For the field, it
would indicate that peripheral visual attenuation on passthrough, the cheapest and most deployable
DR, does not by itself buy back distraction costs, redirecting effort toward the channels the
overlay cannot touch: semantic knowledge of the distractor, auditory leakage, and interruption
timing. For this dissertation, the contributions that survive are exactly the ones constructed to
be result-independent: the design space and its tier constraints (C1), the reference implementation
and its transferable engineering lessons (C2), the benchmark characterisation of a consumer VST
platform (C3), and a pre-registered, falsifiable evaluation method for DR attention claims (C4),
which a null result would demonstrate rather than undermine. A field that currently argues DR's
promise largely from simulation would gain its first hard, hardware-grounded boundary.

\textbf{In either future, the method is part of the contribution.} The evaluation apparatus built for
this dissertation, a single primary endpoint with a smallest effect size of interest, equivalence
margins for safety claims, a gated pilot whose failure aborts rather than degrades the main study,
and a decision table committed before data collection, is not standard practice at
masters-project scale, where underpowered multi-hypothesis designs that can only end inconclusively
remain common. The apparatus costs little (its artefacts are documents and thresholds, not
equipment) and converts every outcome into a claim with stated bounds. Whichever row of the
decision table occurs, the dissertation demonstrates the apparatus by using it in public, which is
a transferable contribution to how small-N XR evaluations can be run.

Both futures are written here with equal seriousness because the study was designed so that either
one is a finding. That property, not any particular row of the decision table, is the standard
this dissertation set out to meet.
